\section{Stable pointed rational curves}
\label{sec:F}

Throughout this section the ground field is $\mathbb C$, the markings are
individually labelled, and $n\geq3$.  We write $X_n=\overline M_{0,n}(\mathbb
C)$ and $U_n=M_{0,n}(\mathbb C)$.  A stratum means a stratum for the stable
labelled dual tree, not a stratum in a chosen triangulation.  We distinguish
the exit category $C_n$ from the entrance category $T_{0,n}=C_n^{\op}$.
Thus arrows of $C_n$ smooth nodes, whereas arrows of $T_{0,n}$ specialize
curves.  All groups below are discrete groups.

\subsection{The four-pointed case}

\begin{statement}{F.1}{Known}{The four-pointed moduli space}
The cross ratio identifies $X_4$ with $\mathbb P^1(\mathbb C)$ and $U_4$
with $\mathbb P^1(\mathbb C)\setminus\{0,1,\infty\}$.  The other three
strata are points, indexed by the partitions $12|34$, $13|24$, and $14|23$.
Writing $G_4=\langle a,b\rangle$ for the free group of rank two, we can
choose positive peripheral elements
\[
 m_0=a,\qquad m_1=b,\qquad m_\infty=(ab)^{-1}.
\]
Each peripheral homomorphism $\mathbb Z\longrightarrow G_4$, $1\mapsto
m_i$, is injective.
\end{statement}
\begin{proof}
An ordered triple of distinct points of $\mathbb P^1$ has a unique
projective transformation to $0,1,\infty$; the fourth point is the cross
ratio.  Stable reduction at the three excluded values gives exactly the
three two-component curves with two markings on each component.  This
also identifies the compactification; see \citet[Catalogue, Facts 1--2,
pp.~551--552]{FKeel1992}.  Removing disjoint small discs around the three
points leaves a pair of pants, which retracts to a wedge of two circles.
The boundary relation, with connecting paths fixed, is
$m_0m_1m_\infty=1$.  Each displayed word is nontrivial in a free group,
which is torsion-free, proving injectivity.
\end{proof}

\begin{statement}{F.2}{Proved}{All the arrows of $C_4$}
A skeleton of $C_4$ has objects $u,p_0,p_1,p_\infty$, automorphism groups
$\Aut(u)=G_4$ and $\Aut(p_i)=1$, and the only noninvertible hom-sets are
\[
 \Hom_{C_4}(p_i,u)=G_4/\langle m_i\rangle.
\]
The quotient consists of left cosets $g\langle m_i\rangle$, formed by
the right action of $\langle m_i\rangle$.  Postcomposition is the left
$G_4$-action, $h\circ[g]=[hg]$.  There are no arrows between distinct
$p_i$, nor from $u$ to a $p_i$.  The opposite of this specified category
is $T_{0,4}$.
\end{statement}
\begin{proof}
The endomorphism space at $u$ is the loop space of the aspherical
three-punctured sphere, hence is equivalent to the discrete group $G_4$.
The endomorphism space at a point stratum is contractible.  The exit
mapping space from $p_i$ to $u$ is the homotopy fibre over $u$ of the
punctured-disc link $S^1\longrightarrow U_4$; radial reparametrization
identifies the usual homotopy-link model with this fibre.  Its components
are the cosets of the image of $\pi_1(S^1)$, and each component is
contractible: the long exact sequence of the homotopy fibre and the
injectivity in F.1 give vanishing of all positive homotopy groups.
Transporting the endpoint around $h\in G_4$ gives the displayed left
action.  The frontier order forbids all remaining noninvertible arrows.
Alternatively the identification with the exit infinity-category follows
by restricting the transporter construction in F.7.  These calculations
specify every hom-space and composition, so they identify the entire
infinity-category, rather than merely its homotopy category.
\end{proof}

\begin{statement}{F.3}{Refuted}{Aspherical strata alone do not imply an
ordinary exit category}
The assertion that every conically stratified space with aspherical
strata has an ordinary exit infinity-category is false.
\end{statement}
\begin{proof}
Stratify $S^2$ by a point $p$ and its complement $\mathbb R^2$.  Both
strata are contractible.  The exit mapping space from $p$ to an interior
point is the homotopy fibre of the link map
$S^1\longrightarrow\mathbb R^2$, hence is $S^1$, and is not discrete.
An ordinary category has discrete mapping spaces.  The nonzero first
homology of the link is also checked from a three-edge simplicial circle
in \texttt{code/F\_checks.py}.  Thus F.2 and F.7 require information about
links and their maps, in addition to asphericity of strata.
\end{proof}

\subsection{The five-pointed case}

\begin{statement}{F.4}{Proved}{The strata and groups for five markings}
There are twenty-six strata in $X_5$: one open stratum $u$, ten
one-dimensional strata $d_I$ indexed by two-element subsets
$I\subset\{1,2,3,4,5\}$, and fifteen point strata $q_{I,J}$ indexed by
unordered pairs of disjoint two-element subsets.  Here $d_I$ represents
the split $I|I^c$, and $q_{I,J}$ has end components carrying $I$ and $J$,
with the remaining marking on the middle component.  Each $d_I$ is a
three-punctured sphere.  The groups of the corresponding exit objects are
\[
 \Aut(u)=G_5=\pi_1\{(x,y)\in(\mathbb C\setminus\{0,1\})^2:x\ne y\},
 \quad \Aut(d_I)=H_I\cong F_2,\quad\Aut(q_{I,J})=1.
\]
An entirely group-theoretic specification of the first group is
\[
 G_5=\ker(B_4\longrightarrow\mathfrak S_4)
       /\langle(\sigma_1\sigma_2\sigma_3)^4\rangle,
\]
where $B_4=\langle\sigma_1,\sigma_2,\sigma_3\mid
\sigma_i\sigma_{i+1}\sigma_i=\sigma_{i+1}\sigma_i\sigma_{i+1},\,
\sigma_1\sigma_3=\sigma_3\sigma_1\rangle$ and the map sends
$\sigma_i$ to $(i,i+1)$.
\end{statement}
\begin{proof}
A stable five-leaf tree has zero, one, or two internal edges.  A single
edge gives a $2|3$ partition, hence ten possibilities.  A trivalent tree
has two disjoint cherries, hence
$\binom52\binom32/2=15$ possibilities.  Gluing the normalized components
identifies the open one-edge stratum with $M_{0,3}\times M_{0,4}$ and
the two-edge stratum with $M_{0,3}^3$.  No automorphism of a labelled
stable genus-zero curve survives, since it fixes three special points on
each normalized component.  Normalizing the first three markings to
$0,1,\infty$ gives the displayed configuration space.  Alternatively,
fix the fifth marking at infinity.  The configuration space of the other
four points in $\mathbb C$, divided by the free affine action, gives
$U_5$.  The affine group is homotopy equivalent to a circle, whose loop
rotates all four points once and represents the central full twist in
$P_4=\ker(B_4\to\mathfrak S_4)$.  Normalization of two of the four
points trivializes the affine bundle, so its homotopy sequence gives
exactly the quotient displayed above.  The asphericity proof in F.6
identifies these fundamental groups with the full homotopy types of the
strata.  The finite enumeration is independently reproduced in
\texttt{code/F\_checks.py}.
\end{proof}

\begin{statement}{F.5}{Known}{A complete specification of $C_5$}
Fix an oriented five-marked sphere $S$.  Let $G=G_5$ act on isotopy
classes of essential multicurves, by orientation-preserving
homeomorphisms fixing every marking.  Choose one multicurve $\sigma_s$
for every stratum $s$ of F.4: $\sigma_u=\varnothing$,
$\sigma_{d_I}=\{c_I\}$, and
$\sigma_q=\{c_{q,I},c_{q,J}\}$.  Put
$\Delta_s=\langle t_c:c\in\sigma_s\rangle\subset G$, where $t_c$ is a
positive Dehn twist.  Then every hom-set is specified by
\[
 \Hom_{C_5}(s,t)=
 \{g\in G:g\sigma_s\supseteq\sigma_t\}/\Delta_s,
 \qquad [h]\circ[g]=[hg].                         \tag{F.1}
\]
In particular
\[
 \Hom(d_I,u)=G/\langle t_{c_I}\rangle,
 \qquad \Hom(q,u)=G/\Delta_q,
 \qquad \Delta_q\cong\mathbb Z^2.
\]
If $I$ is one of the two cherries of $q$, then
\[
 \Hom(q,d_I)\cong H_I/\langle m_{q,I}\rangle,
                                                        \tag{F.2}
\]
where the three $m_{q,I}$, as $q$ ranges over the three points in
$\overline d_I$, are the three peripheral elements of
$H_I\cong\langle a_I,b_I\rangle$, namely $a_I,b_I,(a_Ib_I)^{-1}$
after choices of connecting paths.  These hom-sets are empty for all
other pairs of distinct strata.  Formula (F.1) identifies this ordinary
category with the exit infinity-category of $X_5$.
\end{statement}
\begin{proof}
The identification with exit paths is the case $(g,n)=(0,5)$ of
\citet[Construction 3.18 and Corollary 6.6]{FJansen2023}.  We verify all
the displayed algebraic data.  Cutting $S$ along a multicurve gives the
exact sequence
\[
 1\longrightarrow\Delta_s\longrightarrow
 \operatorname{Stab}_G(\sigma_s)\longrightarrow
 \prod_v\pi_1(M_{0,\operatorname{val}(v)})\longrightarrow1.       \tag{F.3}
\]
Indeed a homeomorphism of the cut, punctured components lifts across the
annuli; the ambiguity consists precisely of twists of those annuli.
Distinct disjoint essential curves have independent twists, as can be
checked by arcs transverse to one curve and disjoint from all the others.
In genus zero with fixed leaf labels the dual tree has no automorphisms,
so there is no additional quotient or permutation group.  Equation
(F.3) gives $\operatorname{Stab}(c_I)/\langle t_{c_I}\rangle=H_I$ and
$\operatorname{Stab}(\sigma_q)=\Delta_q$.

For a fixed incident pair $(q,I)$ choose
$r_{q,I}\in G$ carrying $c_{q,I}$ to $c_I$.  Every transporter from $q$
to $d_I$ is then $k r_{q,I}$ with
$k\in\operatorname{Stab}(c_I)$.  The conjugate
$r_{q,I}\Delta_qr_{q,I}^{-1}$ is generated by $t_{c_I}$ and the twist
about the remaining curve in the four-pointed component.  Modulo
$\langle t_{c_I}\rangle$ the latter is its peripheral element
$m_{q,I}$.  This proves (F.2), including the left action of $H_I$.
The split labels show that no other transporters exist.  Finally the
composition in (F.1) is well-defined: replacing $g$ by $g\delta$, with
$\delta\in\Delta_s$, changes the product on the right by
$\delta$; replacing $h$ by $h\epsilon$, with
$\epsilon\in\Delta_t$, changes $hg$ by right multiplication by
$g^{-1}\epsilon g\in\Delta_s$, because
$g\sigma_s\supseteq\sigma_t$.  Associativity is inherited from $G$.
This also checks every bimodule action in the specification.
\end{proof}

\begin{statement}{F.6}{Proved}{Composition through the one-dimensional
strata, and the asphericity of every genus-zero stratum}
In the coordinates of F.5 the remaining composition is
\[
 \bigl(G/\langle t_{c_I}\rangle\bigr)
 \times\bigl(H_I/\langle m_{q,I}\rangle\bigr)
 \longrightarrow G/\Delta_q,
 \quad ([g],[h])\longmapsto[g\widetilde h r_{q,I}],              \tag{F.4}
\]
where $\widetilde h\in\operatorname{Stab}(c_I)$ lifts $h$.
Together with multiplication and the bimodule actions, this gives all
compositions in $C_5$.  More generally the stratum of a stable labelled
tree $\tau$ is exactly
\[
 U_\tau\cong\prod_{v\in V(\tau)} M_{0,\operatorname{val}(v)}(\mathbb C),
                                                               \tag{F.5}
\]
and each factor is aspherical.
\end{statement}
\begin{proof}
Changing the lift $\widetilde h$ by a power of $t_{c_I}$ changes its
product with $r_{q,I}$ by an element of $\Delta_q$; changing $h$ by a
power of $m_{q,I}$ does the same by the second generator of
$r_{q,I}\Delta_qr_{q,I}^{-1}$.  Changing $g$ by a power of
$t_{c_I}$ is harmless for the same reason.  Thus (F.4) is independent of
representatives and is precisely (F.1).  A nonidentity composable chain
between distinct strata has length at most two, so nothing remains to
specify.  Different routes through $d_I$ and $d_J$ agree exactly when
their products determine the same coset $g\Delta_q$.  We have not
silently chosen simultaneous representatives $c_I$ for which all
compatible pairs are disjoint; the transporters $r_{q,I}$ retain the
necessary change-of-marking information.

For (F.5), normalize the curve and label the inverse images of nodes by
the half-edges of $\tau$.  This defines a map from the stratum to the
displayed product.  Gluing the two labelled points of each edge gives
its inverse.  A leaf-fixing automorphism of a stable tree fixes every
vertex: each vertex is the median of three leaves in distinct incident
branches.  Hence the inverse constructions introduce no finite quotient.
This proves the product assertion, including the absence of stack
isotropy.  It is also the genus-zero specialization of the gluing
description in \citet[Catalogue, Fact 2]{FKeel1992}.

The normalization of three markings identifies $M_{0,k}$ with the
ordered configuration space of $k-3$ points in
$\mathbb C\setminus\{0,1\}$.  Forgetting the last point is a locally
trivial bundle with fibre a finitely punctured plane; the bundle is
\citet[Theorem 3]{FFadellNeuwirth1962}.  For completeness, disjoint small
discs about a configuration and isotopies supported in these discs give
local trivializations by moving the punctures.  The punctured plane
retracts to a graph, and its universal cover is a tree.  Induction on
$k$, starting from $M_{0,3}=*$, and the homotopy exact sequence of this
bundle show that $\pi_i(M_{0,k})=0$ for $i\geq2$.  Products preserve this
property.  Thus the asphericity assertion includes $k=3$ and every
stable tree, without invoking compactified moduli asphericity.
\end{proof}

\subsection{The category for arbitrary numbers of markings}

\begin{statement}{F.7}{Known}{The ordinary exit category and the
Charney--Lee model}
Let $S_{0,n}$ be an oriented $n$-punctured sphere and
$\Gamma_n=\pi_0\operatorname{Homeo}^+(S^2;1,\ldots,n)$, with markings
fixed pointwise.  Let $\mathcal C_n$ be the set of isotopy classes of
finite essential multicurves, including the empty one.  Define a small
ordinary category $C_n$ by
\[
 \operatorname{Ob}(C_n)=\mathcal C_n,\qquad
 \Hom_{C_n}(\sigma,\tau)
 =\{g\in\Gamma_n:g\sigma\supseteq\tau\}/\Delta_\sigma,\qquad
 [h]\circ[g]=[hg].                                      \tag{F.6}
\]
The quotient is by the right action of $\Delta_\sigma$, so its elements
are left cosets $g\Delta_\sigma$.
Then $NC_n\simeq\Exit(X_n)$ as infinity-categories over the poset of
dual trees.  The category is EI and its automorphism group at $\sigma$
is $\prod_v\Gamma_{\operatorname{val}(v)}$.  Isomorphism classes of
objects are stable labelled dual trees.  The category
$T_{0,n}=C_n^{\op}$ is equivalent to the original Charney--Lee category
of stable labelled topological curves and isotopy classes of
degenerations.
\end{statement}
\begin{proof}
The exit equivalence is \citet[Corollary 6.6]{FJansen2023}; the
identification with the original Charney--Lee category is that paper's
Proposition 3.23.  These results establish an equivalence of
infinity-categories, not merely a weak equivalence of classifying
spaces.  Their hypotheses apply because $X_n$ is the analytification of
the stable-curve moduli stack, and in labelled genus zero every
stabilizer of an algebraic curve is trivial.  The boundary is a simple
normal-crossings divisor by \citet[Catalogue, Fact 1]{FKeel1992}.
The definition and the typing of composition have been verified in
F.5; that proof uses no restriction to five markings.  If
$g\sigma\supseteq\sigma$, equality follows from finiteness and equal
cardinality; consequently every endomorphism is invertible.  Equation
(F.3) proves the asserted automorphism group in every valence.
Smoothing annuli identifies the orbits of multicurves with labelled
stable trees, and the converse follows by gluing homeomorphisms of
the normalized components across annuli.  This proves the remaining
algebraic assertions.  A skeleton is finite on objects but generally
infinite on arrows.  We retain the nonskeletal model (F.6) when making
claims about strict pullbacks of functors.
\end{proof}

\begin{statement}{F.8}{Proved}{Coefficients and the necessary opposite}
For a compactly generated presentable infinity-category $\cC$ there
is an equivalence
\[
 \Sh_{\rm cbl}(X_n;\cC)\simeq
 \Fun(C_n,\cC)=\Fun(T_{0,n}^{\op},\cC).              \tag{F.7}
\]
In particular $\HOT(T_{0,n})\simeq\Ho\Sh_{\rm cbl}(X_n;\Spaces)$.
If $\D=\D_{\cC}$ is the presheaf-variance derivator of such a
$\cC$, the same assertion holds at every diagram shape.  For a bare
derivator, $\D(T_{0,n})$ and its shifted values define a theory of
categorical constructible coefficients; no equivalence with a separately
defined category of topological $\D$-valued sheaves follows from Der
1--4 alone.
\end{statement}
\begin{proof}
Equation (F.7) is the coefficient conclusion of
\citet[Theorem 6.7]{FJansen2023}.  Apply it also to the diagram
coefficient category $\Fun(I^{\op},\cC)$, or use the exponential law,
to obtain the assertion at shape $I$.  Taking homotopy categories gives
the displayed derivator statements.  The last sentence separates a
definition from a theorem: an abstract derivator supplies the categories
$\D(A)$ and homotopy Kan extensions, but does not supply sheaves on the
topological site of $X_n$.  Nor does the cited coefficient theorem
assert its conclusion for every presentable $\cC$ without an additional
coefficient hypothesis.  Finally, the prompt's literal combination
``$T_{0,n}=\Exit(X_n)$'' and
``$\HOT(T_{0,n})=$ constructible sheaves'', with the fixed presheaf
variance, requires an opposite.  Our notation makes that correction
explicit rather than changing the direction of exits.
\end{proof}

\subsection{Clutching, forgetting markings, and base change}

\begin{statement}{F.9}{Proved}{The actual functor that forgets a marking}
Filling the puncture $n+1$ induces a homomorphism
$\phi:\Gamma_{n+1}\to\Gamma_n$ and a function
$q:\mathcal C_{n+1}\to\mathcal C_n$ which discards newly peripheral
curves and replaces newly isotopic curves by a single isotopy class.
The assignments
\[
 Q:C_{n+1}\longrightarrow C_n,
 \quad\sigma\longmapsto q\sigma,\quad[g]\longmapsto[\phi(g)]
                                                               \tag{F.8}
\]
define a strict functor.  Its opposite
$F:T_{0,n+1}\to T_{0,n}$ represents the geometric morphism
$f:X_{n+1}\to X_n$ which forgets the last marking and stabilizes.
Repeated forgetting and relabelling are compatible with composition;
on chosen skeletons the compatibility is by specified natural
isomorphisms arising from the chosen markings.
\end{statement}
\begin{proof}
The operation on curves is independent of isotopy, preserves inclusion,
and is equivariant for $\phi$.  A twist becomes the twist about the
image essential curve or the identity if the image is peripheral.
Thus $\phi(\Delta_\sigma)\subseteq\Delta_{q\sigma}$ and the coset map
is well-defined.  If $g\sigma\supseteq\tau$, then
$\phi(g)q\sigma\supseteq q\tau$, which checks the source and target
of every arrow.  The formula respects multiplication, identities, and
composition.  Deleting two punctures successively and reducing the
curve systems gives the same set of surviving essential isotopy classes
as deleting both at once, so the functors on the fixed marked-surface
models compose strictly.

For a permutation of the labels choose an orientation-preserving
homeomorphism $r$ inducing it.  Relabelling acts on $C_n$ by
$\sigma\mapsto r\sigma$ and $[g]\mapsto[rgr^{-1}]$.
If another lift is $r'=ar$ with $a\in\Gamma_n$, the arrows
$[a]:r\sigma\to r'\sigma$ are a natural isomorphism: both composites
in its naturality square are represented by $a r g r^{-1}$.
For two permutations compare their product lift with the chosen lift
of the product by this same formula.  Associativity of multiplication
gives the coherence identity.  Thus no splitting of the full mapping
class group onto the permutation group is assumed.  Filling a relabelled
puncture commutes with these assignments because filling is equivariant
on curve classes and homeomorphisms.

A pinching map of marked surfaces commutes with filling a puncture
followed by contraction of its unstable component: on the complement
of the pinched annuli this is the same continuous map; on an unstable
sphere it is the canonical collapse.  Twisting a pinched annulus is
precisely the equivalence already divided out in (F.6).  Therefore the
induced functor on degeneration categories is (F.8) with its direction
reversed.  Equivalently, the maps on strata and on the normal annuli
give (F.8) under exit paths.  This supplies the comparison with the
geometric map, not just an abstract functor with the correct object map.
\end{proof}

\begin{statement}{F.10}{Proved}{Properness of the specified forgetful
functor}
For the full multicurve models (F.6), the entrance functor
$F:T_{0,n+1}\to T_{0,n}$ is $W_\infty$-proper in Grothendieck's sense.
Thus its right-Kan base-change transformation is invertible for every
strict base-change square and every derivator in the conventions of
this paper.  The assertion also holds for repeated forgetting of
markings.  It is an assertion about these specified functors, not about
arbitrary replacements by equivalent skeletal categories.
\end{statement}
\begin{proof}
We use the outgoing-arrow lifting criterion of Part D, or its
incoming-arrow dual on exit categories.  Fix $\sigma\in\mathcal C_{n+1}$,
put $\tau=q\sigma$, and fix an arrow $\beta:\eta\to\tau$ of $C_n$.
The relevant category has objects arrows
$a:\rho\to\sigma$ in $C_{n+1}$ with $q\rho=\eta$ and
$Q(a)=\beta$; its morphisms are triangles commuting over $\eta$.
If $\beta=[b]$, choose $\gamma\in\Gamma_{n+1}$ with
$\phi(\gamma)=b$.  The assignments
$(\rho,[g])\mapsto(\rho,[\gamma^{-1}g])$ and the identity on
triangle morphisms give an isomorphism to the lifting category with
target $\gamma^{-1}\sigma$ and base arrow
$[1]:\eta\to b^{-1}\tau$; the inverse multiplies by $\gamma$.
Thus, after this specified comparison, we can write the test as
$\beta=[1]:\eta\to\tau$ with $\eta\supseteq\tau$.
Surjectivity is the Birman exact sequence
\[
 1\longrightarrow\pi_1(S_{0,n})\longrightarrow\Gamma_{n+1}
 \xrightarrow{\phi}\Gamma_n\longrightarrow1;                 \tag{F.9}
\]
for $n=3$ it is also immediate from F.1.  For $n\geq4$ we use
\citet[\S2.1]{FKentLeiningerSchleimer2009}.

An object is now $(\rho,[g])$ with
$g\rho\supseteq\sigma$ and $\phi(g)\in\Delta_\eta$.
The arrow $[g]:\rho\to g\rho$ is an isomorphism over $\eta$ and
identifies this object with $(g\rho,[1])$.  The full subcategory of
objects $(\rho,[1])$ is exactly the poset
\[
 P(\eta;\sigma)=
 \{\rho\in\mathcal C_{n+1}:\rho\supseteq\sigma, q\rho=\eta\},
 \quad \rho\leq\rho'\ \Longleftrightarrow\ \rho\supseteq\rho'. \tag{F.10}
\]
Indeed a triangle between two such objects is represented by $h$ with
$[h]=[1]\in\Hom(\rho,\sigma)$, hence $h\in\Delta_\rho$.
It therefore exists precisely when $\rho\supseteq\rho'$ and is then
unique.  This proves an equivalence of categories with (F.10), including
its arrows, rather than only a bijection on connected components.

We prove that (F.10) has contractible nerve.  First suppose
$\sigma=\varnothing$.  Realize $\eta$ by disjoint geodesics in a
complete hyperbolic metric on the $n$-punctured sphere.  In its universal
disc their preimages are disjoint proper lines.  The graph dual to
these lines is a tree: it is connected, and every line separates the
disc, so the corresponding edge separates the graph.  The vertices
are components of the complement, and the edges are the lifted lines.
The deck group $K=\pi_1(S_{0,n})$ acts with vertex stabilizers
$\pi_1(V)$ for the components $V$ of $S_{0,n}\setminus\eta$ and edge
stabilizers the cyclic subgroups generated by the curves of $\eta$.

Point-pushing identifies lifts of $\eta$ having the same number of
curves with the vertices of this tree.  In detail, moving the forgotten
puncture into a component $V$ gives a lift; two moves give the same
isotopy class exactly when their difference is a loop in $V$, so these
lifts are the cosets $K/\pi_1(V)$.  A lift with one extra curve whose
two images are isotopic places the forgotten puncture between two
parallel curves.  Such lifts are the cosets of the corresponding edge
subgroup and are incident to the two lifts obtained by deleting one of
the parallel curves.  This is precisely the barycentric incidence
description of the dual tree.  This description is also the content
of \citet[Theorem 7.1]{FKentLeiningerSchleimer2009}; their Lemma 2.1
checks that at most one duplication occurs.

There is one additional possibility because (F.10) allows curves that
become peripheral: a disc contains the forgotten puncture and exactly
one retained puncture.  Such a lift adds one peripheral leaf to the
regular lift in the adjacent component.  Its stabilizer under
point-pushing is the cyclic group of that puncture.  It has a unique
proper face still mapping to $\eta$, obtained by deleting this extra
curve.  There cannot be two such peripheral curves, nor a peripheral
curve and a duplicated pair, because only one puncture is filled.
Thus the nerve is a subdivided tree with leaves attached.  It is
contractible.  When $\eta=\varnothing$ there is a central vertex
$\varnothing$ and only these leaves, so the same conclusion holds
without using a nonempty curve complex.

For general $\sigma$, cut the pointed sphere along $\sigma$ before
filling the puncture.  Every component except the one containing the
forgotten puncture has a unique lift of each prescribed further
refinement in $\eta\setminus\tau$: isotopy classes there are unchanged
by filling a puncture on another component.  If the distinguished
component remains stable after filling, the remaining choices are
exactly the preceding tree-and-leaves calculation on that component,
whose other special points are treated as punctures.  If it becomes
unstable, it had exactly three special points before filling.  They
are either the forgotten marking, a retained marking, and a node, or
the forgotten marking and two nodes.  In the first case it is a fixed
peripheral leaf; in the second it is a fixed inserted edge.  Preserving
$\sigma$ fixes this placement, and the additional curves of $\eta$
lift uniquely in the other components.  In either case (F.10) is a
singleton.  These cases exhaust the possibilities.  We have proved
contractibility of every required lifting category, and hence
properness.  Proper functors are closed under composition, which proves
the final assertion.
\end{proof}

\begin{statement}{F.11}{Refuted}{Forgetting a marking need not be smooth}
The proper entrance functor $T_{0,5}\to T_{0,4}$ is not
$W_\infty$-smooth.  The first stable forgetful map
$T_{0,4}\to T_{0,3}=e$ is both proper and smooth, so five markings are
the smallest instance of this failure in the genus-zero forgetful
sequence.
\end{statement}
\begin{proof}
Take the upstairs point stratum $x=q_{12,34}$, with marking $5$ on the
middle three-pointed component.  Its two curves $c_1,c_2$ become the
same essential curve $c$ after forgetting $5$.  Write
$\Delta_x=\langle t_1,t_2\rangle\cong\mathbb Z^2$ and
$\phi(t_1)=\phi(t_2)=t_c$.  In exit variance consider the arrow
$[1]:q x=\{c\}\to\varnothing$.  Its outgoing lifting category is the
test for smoothness in entrance variance.

The target of any such lift has empty image under $q$.  It is therefore
either the empty curve system or a single curve enclosing $5$ and one
retained marking.  The latter cannot be a subset of any translate of
$\sigma_x$: its split label is different from both $12|345$ and
$34|125$, and pure mapping classes preserve split labels.  Hence the
only possible target is the open object.  Put
$K=\ker(\phi)$.  Objects of the lifting category are
$\{g\Delta_x:\phi(g)\in\langle t_c\rangle\}$, and arrows between
them are postcomposition by elements of $K$.  Because
$\phi(\Delta_x)=\langle t_c\rangle$, this is the transitive action
groupoid of $K$ on $K/(K\cap\Delta_x)$.  Its stabilizer is
\[
 K\cap\Delta_x=\langle t_1t_2^{-1}\rangle\cong\mathbb Z.
\]
The lifting category therefore has classifying space $B\mathbb Z$,
not a contractible space.  This verifies every hypothesis and the
failure of the smoothness criterion.  Independence of $t_1,t_2$
follows from their disjoint essential supports as in F.5; in particular
the stabilizer is not killed by the quotient defining the arrows.
The check of the exponent map $(r,s)\mapsto r+s$ is included in the
code, although the counterexample category itself is infinite.  A
functor to the terminal category passes both lifting tests because its
tests are ordinary slices with an initial or terminal object.  This
settles the minimality assertion within the sequence.
\end{proof}

\begin{statement}{F.12}{Proved}{Clutching as a closed categorical
embedding}
Let $L=A\amalg B$ be a partition with $|A|,|B|\geq2$.  Clutching the
extra marked points gives an isomorphism
\[
 \overline M_{0,A\cup\{*\}}\times
 \overline M_{0,B\cup\{*\}}\xrightarrow{\sim}D_{A|B}\subset X_{|L|}
                                                               \tag{F.11}
\]
onto the closed boundary divisor.  In entrance variance the induced
functor can be represented by the fully faithful cosieve on the
objects whose dual trees contain the split $A|B$.  It is therefore
$W_\infty$-proper.  On exit categories the corresponding inclusion is
a fully faithful sieve.  The product description and this inclusion
are compatible with clutching more than one node.  The inclusion of
the fixed-curve product model into the full multicurve model is not
itself $W_\infty$-proper; the proper representative is its replete
enlargement just specified.
\end{statement}
\begin{proof}
Normalization at the distinguished node inverts (F.11); fixed labels
distinguish the two branches, including when their cardinalities are
equal.  This is \citet[Catalogue, Fact 2]{FKeel1992}.
Choose a fixed curve $c$ separating $A$ from $B$, and first use only
multicurves containing $c$.  An arrow between two such objects is
represented by a mapping class preserving $c$: in the image of the
source multicurve there is exactly one curve with the labelled split
$A|B$, so it must be the fixed target curve.  The stabilizer of $c$
modulo its twist is the product of the mapping class groups of the
two punctured components.  This proves the product identification of
categories and their full faithfulness.

For the actual cosieve take the full replete subcategory of all
multicurves containing some curve with split $A|B$.  It is equivalent
to the fixed-$c$ product just described, and it is a sieve in exit
variance since refinement retains the split.  Its opposite is the
required literal cosieve.  This distinction matters: the inclusion
of the fixed-$c$ subcategory alone is not replete in the nonskeletal
model.  In fact choose a pure mapping class taking $c$ to a distinct
curve $c'$; such a class exists by twisting about an essential curve
intersecting $c$, since both sides of $c$ contain at least two markings.
The entrance isomorphism $\{c\}\to\{c'\}$ has no lift under the
fixed-$c$ inclusion: its target object $\{c'\}$ is not in that
inclusion's literal image.  Its outgoing-arrow lifting category is
empty, proving the claimed failure of strict properness for that
particular representative.  Repeating the normalization argument at
several disjoint nodes gives the stated compatibility.  The properness
of the replete full cosieve follows from the lifting criterion of Part D.
\end{proof}

\begin{statement}{F.13}{Proved}{Strict categorical fibres and geometric
fibres of forgetting}
For the full model (F.8), the strict fibre $F^{-1}(\tau)$ is equivalent
to the entrance category of the stable $n$-pointed curve represented by
$\tau$, stratified by its smooth component complements, its marked
points, and its nodes.  This is the topological fibre of
$X_{n+1}\to X_n$.  In particular the fibre of $X_5\to X_4$ is $S^2$
over $U_4$ and $S^2\vee S^2$ over any point stratum; the map is not a
topological fibre bundle and is not a submersion of the underlying
smooth manifolds.
\end{statement}
\begin{proof}
An object upstairs that stabilizes to a fixed nodal curve places the
extra point either on the smooth locus of a component, at a retained
marking by adding a three-pointed bubble, or at a node by inserting a
three-pointed bridge.  These are exactly the strata of that curve.
In the strict categorical fibre, two placements of the same type
are isomorphic by point-pushing.  At a smooth-component placement the
automorphism group is the kernel of forgetting in its component
mapping class group, namely the fundamental group of that component
with its special points removed, by (F.9).  At a bubble or bridge the
automorphism group is trivial.  The incidence bimodule at each adjacent
marked point or node is the set of left cosets of the corresponding
peripheral cyclic subgroup in that component group, by the same
annular calculation as F.2.  There are no chains through three
different strata.  This identifies every hom-set, action, and
composition with the entrance category of the curve, and proves the
claimed categorical equivalence.

The geometric assertion about the fibre is the universal-curve
interpretation of forgetting and stabilizing; it is also described in
\citet[Introduction, pp.~546--548]{FKeel1992}.  A smooth four-pointed
curve is one projective line.  A boundary four-pointed curve is two
projective lines meeting in one point.  Their second integral homology
has rank one and two respectively.  Fibres of a topological fibre
bundle over a connected base are homeomorphic and have equal
homology.  A proper smooth submersion would be a fibre bundle, and
in this instance the singular fibre also directly excludes a
submersion.  This geometric argument is independent of, and weaker
than, the categorical obstruction in F.11.
\end{proof}

\begin{statement}{F.14}{Proved}{The coefficient functors and their
base-change properties}
For a compactly generated presentable coefficient infinity-category,
pullback along clutching or forgetting corresponds under (F.7) to
restriction along the entrance functors above.  Their ordinary
topological direct images preserve constructibility and correspond
to right Kan extension.  For every derivator, the specified proper
categorical models satisfy right-Kan base change for all strict
category pullbacks.  Geometric proper base change holds for the
analytic maps and their pullbacks among locally compact Hausdorff
spaces.  The two assertions agree
whenever the strict pullback category is identified with the entrance
category of the geometric pullback, in particular for the point
fibres in F.13 and for the closed-boundary squares in F.12.  There is
no smooth-base-change assertion for forgetting: F.11 refutes that
hypothesis already at five markings.
\end{statement}
\begin{proof}
The geometric maps preserve strata and send exit paths to exit paths,
so the pullback assertion follows from exodromy and F.9--F.12.  For a
closed inclusion, direct image is extension from the closed subset,
and its values and gluing maps are the corresponding right Kan
extension; the local open--closed calculation is the one in Part D.
Forgetting is a proper algebraic morphism.  Its restriction over a
fixed target stratum is locally topologically trivial as a family of
stratified nodal curves: choose disjoint coordinate discs around the
labelled special points on each normalized component, move the
points by isotopies in these discs, and glue the fixed node pairs.
Shrink to a contractible coordinate ball $B$ in the target stratum,
and write its stratified trivialization as $B\times C$, with $C$ the
stratified nodal fibre.  By the definition of exit simplices, the exit
category of this product is
$\Pi_\infty(B)\times\Exit(C)\simeq\Exit(C)$.  Applying coefficient
exodromy to the product shows that the constructible sheaf on it is
pulled back from $C$, up to an equivalence compatible with restriction
to each fibre.  This proves the required locally constant transport;
local triviality of the underlying spaces alone would not supply that
conclusion for arbitrary sheaves.  Compactness of the fibre and proper
base change identify
the stalk of the direct image with sections on this fibre.  They
therefore make the direct image locally constant on the target
stratum.  This proves preservation of constructibility.

Since both inverse image and direct image preserve the respective
constructible subcategories, their adjunction restricts to those
subcategories.  Under (F.7), inverse image is restriction.  Uniqueness
of its right adjoint therefore identifies the restricted direct image
with right Kan extension, including its unit and counit.  In addition,
F.10 computes right Kan extension at $\tau$ as the homotopy limit over
the strict fibre, and F.13 identifies that limit with sections of the
restriction to the geometric fibre.  These are the specified fibrewise
base-change comparisons.
For spaces, geometric proper base change is
\citet[Corollary 7.3.1.18]{FLurieHTT2009}; testing with compact generators
extends it to the indicated coefficient categories, because mapping
out of a compact generator commutes with the filtered stalk colimit
and with the limits in the direct image.  The categorical assertion
uses only the properness theorem of Part D and hence applies also to
bare derivators.  For a square obtained by pulling back a replete
boundary inclusion, the strict pullback is the full subcategory on
objects whose image lies in that boundary.  Its geometric counterpart
is the closed union of strata with exactly this condition.  The exit
category of a closed union of strata is the full subcategory on those
strata: along an exit simplex all strata precede its last stratum, and
closedness retains all these predecessors.  This proves the asserted
comparison for boundary squares.  Nothing in this argument identifies an arbitrary
strict category pullback with a geometric pullback.  That comparison
is extra data, and its verified instances were specified in the
statement.  Likewise the left adjoint of constructible pullback is a
left Kan extension; it should not be confused with the topological
exceptional direct image, whose adjunction is with exceptional
inverse image.
\end{proof}

\begin{statement}{F.15}{Proved}{The explicit map from five to four
markings}
Forget marking $5$.  The open object maps to the open object.
A divisor object $d_I$ maps to the open object when $5\in I$ and to
the boundary partition $I|\{1,2,3,4\}\setminus I$ otherwise.
Every $q_{I,J}$ maps to the partition defined by its surviving
essential split; if $5$ is the middle marking then both splits
survive as the same partition.  On all hom-sets the map is the coset
map $[g]\mapsto[\phi(g)]$ of F.9.  In particular a twist about a
curve enclosing $5$ and one retained mark maps to $1$, and all other
single-curve twists map to the twist about their image split.
\end{statement}
\begin{proof}
Deleting $5$ makes the side $\{5,i\}$ unstable and removes its edge.
Every other divisor split remains stable.  For a two-cherry tree,
deleting an endpoint marking contracts its unstable endpoint
component and leaves the other split; deleting the middle marking
contracts that component and identifies its two incident edges.
These are all the cases.  Annular twists give the arrow formulas,
which are independent of chosen representatives by F.9.  Over each
boundary partition there are two divisor strata and five point
strata, and over the open stratum there are the open stratum and
four divisor strata.  The program enumerates these fibres and all
incidence relations, and checks that the object map is monotone.
The twenty-six-object description, the explicitly defined group
$G_5$, the coset bimodules (F.1)--(F.2), and the multiplication (F.4)
thus give a complete specification of $T_{0,5}\to T_{0,4}$ after
opposites and choices of skeleton markings.
\end{proof}

\subsection{What the Teichm\"uller tower comparison says}

\begin{statement}{F.16}{Known}{The relevant definitions in Grothendieck's
texts}
In \emph{La Longue Marche}, \S19, the notation is
$A_g=\Aut(X_g)$ for the topological homeomorphism group of a compact
orientable genus-$g$ surface,
$A_{g,\nu}=\{u:u(S_{g,\nu})=S_{g,\nu}\}$,
$A_{g,\nu}^{!}=\{u:u|_{S_{g,\nu}}=\id\}$, and
\[
 \Gamma_{g,\nu}=\pi_0 A_{g,\nu},\qquad
 \Gamma_{g,\nu}^{!}=\pi_0 A_{g,\nu}^{!}.               \tag{F.12}
\]
Our $\Gamma_n$ is the orientation-preserving subgroup of
$\Gamma_{0,n}^{!}$, not the full group allowing permutations or
orientation reversal.  The modular topos $M_{g,\nu}$ of \S26 is the
modular topos of smooth curves of type $(g,\nu)$; its complex
topological fundamental group is the corresponding oriented
Teichm\"uller group.  It is not the constructible-sheaf category of
the compactification.  The operation called ``bouchage de trous'' in
\S27 induces the forgetful operation of F.9 on the open stratum.
\end{statement}
\begin{proof}
The definitions (F.12) are the literal equations (3)--(11) in the
attached \S19 transcription, manuscript p.~102, scans 204--205.
The superscript $!$ imposes fixed labels; orientation preservation is
an additional restriction, not part of the symbol $A_g$ in those
equations.  This resolves the potential mismatch between analytic
moduli with labelled markings and the whole homeomorphism group.
The supplied P034 extract starts with the continuation of \S26;
for its preceding definition we opened the Grothendieck Circle
transcription, Part Ib, \S26, equations (19)--(25), printed pp.~5--6.
There $M_{g,\nu,\mathbb C}$ is identified through its Teichm\"uller
universal cover, and equation (24) is the universal-curve fundamental
group extension.  The topological restriction of that extension is
(F.9) with the same label and orientation restrictions.  At \S27,
manuscript pp.~172--173, scans 340--342, filling the punctures in
$I\setminus I'$ is explicitly the operation considered.  On the
surface fundamental group it kills the normal closure of those
peripheral loops, by van Kampen.  On homeomorphism components it is
the homomorphism $\phi$ of F.9.  This proves the asserted
compatibility while making no assertion about the arithmetic
extensions discussed on the surrounding pages
\citep[\S\S19,26--27]{FGrothendieckLM1981}.
\end{proof}

\begin{statement}{F.17}{Proved}{The precise comparison with the Lego
groupoid}
The full subgroupoid of $T_{0,n}$ on the open stratum is equivalent
to $B\Gamma_n$.  More generally its maximal subgroupoid is
\[
 \coprod_{\tau} B\!\left(\prod_v\Gamma_{\operatorname{val}(v)}\right).
                                                               \tag{F.13}
\]
The usual topological Teichm\"uller groupoid of $U_n$ is equivalent
to the first of these groupoids.  The equivariant path groupoid formed
from the marking complex of the Lego--Teichm\"uller game and its mapping
class group action presents this groupoid after imposing the same label
and orientation conventions and capping marked boundary circles to
punctures.  The nonequivariant path groupoid of that complex is
contractible.  Neither construction identifies the whole degeneration
category $T_{0,n}$ with a groupoid.
\end{statement}
\begin{proof}
The first two assertions are (F.3) and (F.6): an isomorphism cannot
change the number or labelled configuration of nodes.  For the
opposite category we identify the opposite of each automorphism group
with that group by inversion.  Since $U_n$
is connected and aspherical, its fundamental groupoid is equivalent
to its one-object groupoid $B\Gamma_n$, with the equivalence depending
on a basepoint and connecting paths.  Grothendieck's
\emph{Esquisse}, \S2, explicitly treats the system of open moduli
spaces linked by filling and gluing, and advocates several basepoints
for its fundamental groupoids; see the attached edition,
pp.~9--11 \citep{FGrothendieckEsquisse1984}.

For a genus-zero extended surface $\Sigma$ with $n\geq3$ marked
boundary circles, let $\mathcal L(\Sigma)$ be the
two-complex of markings and elementary moves of
\citet[Main Theorem 4.9]{FBakalovKirillov2000}.  That theorem states
that it is connected and simply connected.  Form the groupoid with
objects its markings and morphisms $M\to M'$ the pairs $(g,[p])$
consisting of a mapping class $g$ and a path $p:gM\to M'$ modulo
homotopy relative to endpoints.  For $q:hM'\to M''$, composition is
$(h,[q])\circ(g,[p])=(hg,[q*(hp)])$, where $hp$ is the translate of
$p$ and the concatenation first traverses $hp$ and then $q$.
Simple connectedness makes the path class unique whenever its endpoints
are fixed.  Consequently the functor sending $(g,[p])$ to $g$ is an
equivalence to the mapping class group's one-object groupoid.  Without
the group action there is exactly one path class between any pair of
markings, so the path groupoid is equivalent to the terminal category.
This distinction is consistent with
\citet[Remark 4.3]{FBakalovKirillov2000}: an elementary marking move
is not itself a homeomorphism of surfaces.
Passing from marked boundary circles to labelled punctures kills the
boundary twists:
extension over the punctured discs gives a surjection, and the
kernel consists exactly of annular twists, by the same cutting
argument as (F.3).  Restricting to fixed labels gives $B\Gamma_n$.
This proves the asserted relation to the Lego presentation.

For $n\geq4$, $T_{0,n}$ has noninvertible degeneration arrows, for
example $u\to p_0$ in $T_{0,4}$.  Such an arrow cannot become
invertible under an equivalence of categories.  Thus the whole
category is not a groupoid.  Its infinity-groupoid completion has
the homotopy type of $X_n$, by F.7 and the groupoid-completion
property of exit paths.  Already $X_4\simeq S^2$ differs from
$U_4\simeq BF_2$.  Ordinary groupoid completion would lose even the
second homotopy group of $S^2$.  Hence neither completion should be
silently substituted for the open Teichm\"uller groupoid.
The input for the separate profinite package is the specified
category (F.6), its opposite convention, the functors F.9 and F.12,
and the properness and nonsmoothness in F.10--F.14.  No profinite
completion or arithmetic action is asserted here.
\end{proof}
